
This appendix formalizes the key mathematical properties that underpin Horizon's economic guarantees.


\section{PT Price Convergence}

\textbf{Theorem A.1.}
Assuming a bounded implied rate $r < \infty$, the PT price converges to its redemption value as expiry approaches:
\[
\lim_{\tau \to 0} P_{PT} = 1.
\]

\textbf{Proof.}

The PT price is defined as:
\[
P_{PT} = e^{-r \cdot \tau_y},
\quad
\text{where}
\quad
\tau_y = \frac{\tau}{Y}
\]
and $Y$ is the number of seconds per year (31,536,000).

As $\tau \to 0$, we have $\tau_y \to 0$. Since $r$ is bounded by construction ($0 \le r \le r_{max}$), the exponent satisfies:
\[
\lim_{\tau \to 0} (-r \cdot \tau_y) = 0.
\]

By continuity of the exponential function:
\[
\lim_{\tau \to 0} e^{-r \cdot \tau_y} = e^{0} = 1.
\]

Therefore, $P_{PT} \to 1$ as $\tau \to 0$. \qed


\section{Conservation Invariant}

\textbf{Theorem A.2.}
For all times $t < T$ (before expiry), the following invariant holds:
\[
1 \, \text{PT} + 1 \, \text{YT} \;\longleftrightarrow\; 1 \, \text{SY}.
\]

\textbf{Proof.}

Before expiry, the protocol exposes two mutually inverse state transitions:

\begin{itemize}
\item \textbf{Split:} Depositing $x$ SY mints exactly $x$ PT and $x$ YT.
\item \textbf{Recombine:} Burning $x$ PT and $x$ YT releases exactly $x$ SY.
\end{itemize}

By construction:
\[
\text{PT.supply} = \text{YT.supply}
\]
and all outstanding PT and YT are fully backed by SY held in the YT contract:
\[
\text{SY}_{\text{YT}} = \text{PT.supply} = \text{YT.supply}.
\]

No operation allows PT or YT to be minted or burned independently. Therefore, for all valid states before expiry:
\[
1 \, \text{PT} + 1 \, \text{YT} = 1 \, \text{SY}.
\]

\qed


\section{No-Arbitrage Pricing Condition}

\textbf{Theorem A.3.}
In the absence of fees, slippage, and execution costs, any deviation from
\[
P_{PT} + P_{YT} = 1
\]
admits a risk-free arbitrage opportunity.

\textbf{Assumptions.}
\begin{itemize}
\item Zero swap fees
\item Zero slippage (infinite liquidity)
\item Zero gas/execution costs
\item Instantaneous execution
\end{itemize}

\textbf{Proof.}

\textbf{Case 1:} $P_{PT} + P_{YT} > 1$

\begin{enumerate}
\item Purchase 1 SY at unit price.
\item Split into 1 PT and 1 YT.
\item Sell PT for $P_{PT}$ and YT for $P_{YT}$.
\end{enumerate}

Net profit:
\[
P_{PT} + P_{YT} - 1 > 0.
\]

\textbf{Case 2:} $P_{PT} + P_{YT} < 1$

\begin{enumerate}
\item Purchase 1 PT for $P_{PT}$ and 1 YT for $P_{YT}$.
\item Recombine into 1 SY.
\item Sell SY at unit price.
\end{enumerate}

Net profit:
\[
1 - P_{PT} - P_{YT} > 0.
\]

In both cases, arbitrageurs are incentivized to trade until the equality
\[
P_{PT} + P_{YT} = 1
\]
is restored, up to execution costs.

\qed

\paragraph{Remark.}
In practice, fees and slippage introduce a bounded arbitrage band around the equality. If total round-trip costs are $\epsilon$, then the no-arbitrage condition relaxes to:
\[
|P_{PT} + P_{YT} - 1| \le \epsilon.
\]
Within this band, arbitrage is unprofitable and deviations may persist.

